% !TEX root = ../print.tex
% !TEX program = xelatex

\begin{definition}[\(L\): rate-block budget]
  \label{def:dsfs_query_error}
  Let \(\SpongeRate>0\) be the sponge rate, \(\SaltLength\) the salt length in alphabet symbols,
  and \(\IPProverPositions,\IPVerifierPositions\) the prover- and verifier-message positions.
  Write \(\IPProverMessageEncodedLength(i),\IPVerifierMessageEncodedLength(i)\) for their encoded lengths.
  Following~\cite[Eqs.~(7)--(8)]{CO25}, define
  \begin{equation*}
  \begin{gathered}
    \SpongeProverMessageBlockCount(i):=\left\lceil\frac{\IPProverMessageEncodedLength(i)}\SpongeRate\right\rceil\ (i\in \IPProverPositions),\qquad
    \SpongeVerifierMessageBlockCount(i):=\left\lceil\frac{\IPVerifierMessageEncodedLength(i)}\SpongeRate\right\rceil\ (i\in \IPVerifierPositions),\\
    L_{\SaltLength}:=\left\lceil\frac\SaltLength \SpongeRate\right\rceil,\qquad
    \SpongeProverMessageBlockCount:=\sum_{i\in \IPProverPositions}\SpongeProverMessageBlockCount(i),\qquad
    \SpongeVerifierMessageBlockCount:=\sum_{i\in \IPVerifierPositions}\SpongeVerifierMessageBlockCount(i),\\
    L:=L_{\SaltLength}+\SpongeProverMessageBlockCount+\SpongeVerifierMessageBlockCount.
  \end{gathered}
  \end{equation*}
\lean{ProtocolSpec.L_totalRateBlocks}
\leanok
\end{definition}

The bound of \(L\) permutation queries for the honest argument verifier is used in the
security proof. The salt and first prover message may share a rate block, so the actual
query count can be smaller.
